\chapter*{Introduction générale}
\addcontentsline{toc}{chapter}{Introduction générale}
\markboth{Introduction générale}{Introduction générale}

L'évolution technologique a profondément transformé plusieurs secteurs d'activité en facilitant l'accès à l'information, l'automatisation des processus et l'amélioration de la communication. Le secteur hôtelier fait partie des domaines concernés par cette transformation, notamment à travers la digitalisation des services destinés aux clients et au personnel. En effet, de nombreux établissements hôteliers utilisent encore des méthodes manuelles pour le check-in, la gestion des réclamations, la communication interne ou le suivi des interventions, ce qui peut entraîner des retards, des erreurs et une perte de traçabilité.

\textbf{Loomens} est une agence de développement web et de transformation digitale située au Cyberparc de Djerba. Elle est spécialisée dans la digitalisation des marques ainsi que dans la création de solutions web et mobiles adaptées aux besoins des entreprises. Dans ce cadre, notre projet consiste à concevoir et développer une plateforme intelligente de services numériques pour les établissements hôteliers, intitulée \textbf{LoomStay}.

La solution \textbf{LoomStay} vise à digitaliser plusieurs services hôteliers à travers une plateforme web, mobile et intelligente. Elle permet aux clients d'accéder à un espace chambre via QR code, d'effectuer un check-in numérique, de consulter les informations de l'hôtel, d'envoyer des messages à la réception et de soumettre des réclamations. Elle offre également aux responsables de maintenance et de housekeeping la possibilité de gérer les plannings, d'assigner les tâches aux employés et de suivre les interventions. Les employés peuvent consulter et valider leurs tâches à l'aide d'une application mobile Flutter et d'un système de scan QR à l'entrée et à la sortie des chambres. Le système intègre des modules d'intelligence artificielle pour la traduction automatique des messages multilingues et la classification des réclamations.

Pour organiser les tâches de développement, nous avons adopté la méthodologie agile Scrum. Cette approche permet de diviser le projet en releases et en sprints, tout en assurant une meilleure adaptation aux changements et une validation progressive des fonctionnalités.

Ce rapport est structuré comme suit :
\begin{itemize}
    \item \textbf{Chapitre 1 --- Contexte général du projet} présente l'organisme d'accueil, le cadre du projet, l'étude de l'existant ainsi que la solution proposée.
    \item \textbf{Chapitre 2 --- Sprint 0 : Analyse et préparation} détaille l'identification des acteurs, les besoins, le backlog, la planification, l'environnement de travail et l'architecture.
    \item \textbf{Chapitre 3 --- Release 1 : Socle système et gestion hôtelière} présente les sprints liés à l'authentification, aux chambres, réservations, QR codes et check-in numérique.
    \item \textbf{Chapitre 4 --- Release 2 : Parcours client et services numériques} décrit l'espace client PWA, les services hôteliers, les événements, la conversion de devises et la messagerie avec traduction automatique.
    \item \textbf{Chapitre 5 --- Release 3 : Réclamations, interventions et intelligence intégrée} présente la gestion des réclamations avec classification IA, l'application Flutter des agents terrain, la gestion des shifts et le suivi des tâches housekeeping.
    \item \textbf{Chapitre 6 --- Build, tests et validation} regroupe les validations techniques réalisées durant le projet.
    \item \textbf{Chapitre 7 --- Déploiement sur serveur privé hôtelier} présente l'architecture de déploiement cible dans un environnement hôtelier privé.
\end{itemize}

Enfin, notre rapport se termine par une conclusion générale qui synthétise le travail réalisé, les apports de la solution développée ainsi que les perspectives d'amélioration futures.
